\section{Discussion}\label{sec:discussion}

\subsection{Why the method behaves as it does}

The ideal-model solver equalizes a field that the device does not produce (Proposition~\ref{prop:floor}), and the refinement equalizes the field that is scored. The ablation places the gain in two components that cannot be separated from each other: the device model in the loop and the conjugate-gradient steps through it. With the model removed the same steps leave the spread at the baseline, \ablBNoModelSig{} percent against \sigmaMsevenThirty{} percent, and without the steps the safeguard returns the warm start, \ablBNoStepSig{} percent. The gain is therefore neither the extra iterations nor the optimizer: it is the optimizer applied to the model of the device that is scored.

The loss is not a third component. Proposition~\ref{prop:losses} shows that for equal traps the efficiency-level loss and the normalized-error loss of Schroff et al.\ are the same function at the point where an outer iteration sets its target, with gradients that differ by a push along the gradient of the total light, weighted by the loss itself. The measurement agrees: the two losses reach spreads that are within each other's seed ranges on the crosstalk device and on the full device, at the same efficiency and the same time. What the efficiency-level loss adds is a statement of what it does with the light, a target at the current efficiency, and a gradient written without the derivative of the normalization; what the normalized-error loss adds is the invariance to the total power, which the safeguard then has to guard. Neither is a source of the gain on this task.

The weight update fails for a reason that the loop explains, though the data do not prove it. The update is an integrator (Proposition~\ref{prop:weights}) and the plant that it controls, a hologram that moves by a few conjugate-gradient steps, can lag the integrator; the outer iterations then overshoot and the spread rises. The ablation measures the consequence on the acceptance task and on the full device, where the update raises the spread several times, and the geometry table shows that on the lattices of the benchmark the update sometimes helps (the ring) and sometimes hurts (the square lattice), so the data support dropping it as a default and keeping it as an option.

\subsection{Failure cases and negative results}

Five results are negative or null and they bound what the paper claims. First, the factor of two over phase-fixed WGS that the lattice behind a crosstalk device shows by two orders of magnitude does not carry to the benchmark's device on the square lattice at 100 iterations, and at the equal wall-clock time of one second the refinement ties phase-fixed WGS on the quasi-periodic set and on both large sets (Section~\ref{sec:results}); the success on the lattice is specific to the device model and the geometry. Second, the normalized-error loss of Schroff et al.\ matches the efficiency-level loss, and on the full device its median is lower, so the proposed loss is not better than the published one; the efficiency-level loss stays the default of the engine, and the paper claims no difference in its favour. Third, the weight update, the part of weighted GS that the earlier design kept, raises the spread by a large factor on the acceptance task and is mixed on the benchmark lattices. Fourth, the cost in time is minutes to tens of minutes per hologram against seconds for the baseline, so the refinement cannot serve a loop that needs a hologram in seconds. Fifth, the benchmark's 100-iteration cells hold two seeds, so no factor on them is a win.

A sixth observation is a measurement without an explanation. On the quasi-periodic set the three forms of the refinement return the same spread on most seeds, and their median equals the median of the weight update; the spread there is of the order of the baseline's. We report it as measured. A reading that the optimizer finds no admissible iterate on that geometry would need the returned iteration, which the table does not hold.

\subsection{Relation to the published methods}

Two published results frame what the paper adds. Schroff et al.\ put the crosstalk model inside conjugate-gradient minimization, lowered the measured error of a large disc, paid about 20 percent of the efficiency and ran about six times longer~\cite{schroff}. The refinement is the point-trap counterpart of that idea in our engine, run on the same device model, and the comparison that the source does not contain, against weighted GS on the same task, shows a factor of \ratioOneGm{} on the lattice behind a crosstalk device at an efficiency ratio of at least \effRatioOneMin{}. That factor is the effect of the device model in the loop, which the source already uses, and it is not an advance of the refinement's own elements: through the same device model and optimizer the source's loss does as well. The two tasks, devices and metrics differ from theirs, so the numbers are not comparable: they measure the disc error by camera, we measure the spread of point traps in simulation. The efficiency cost that they report, about 20 percent, is not observed here on point traps, where the efficiency ratio is at least \effRatioOneMin{}; the two tasks differ, and the paper does not explain the difference.

The phase-fixed rule of Kim et al.~\cite{kim2019} remains the right start and the right baseline: it is the reference of every factor in this paper, and the warm start of the refinement uses its iteration count. Camera feedback is the other published route to the same end and the two are complements: the device model shortens what the camera has to correct, an expectation that is not tested here, and the camera corrects what the model does not contain, such as the flicker, the aberrations and the error of the fitted kernel~\cite{nogrette,schroff}. For the work of Persson et al.~\cite{persson} the paper has only the abstract: a crosstalk-aware modification of a hologram generating algorithm with a spot uniformity improved by more than 100 percent on a device with large crosstalk. Whether that algorithm already contains the combination measured here cannot be said from the abstract, and the paper does not claim that it does not.

\subsection{What the results imply for use}

For a static array behind a device whose crosstalk is modelled, the part to adopt is the device model in a conjugate-gradient step, which is the approach of Schroff et al.; the paper finds no reason to prefer the refinement's own elements to their loss, so either loss serves, the weight update is left off and about 10 outer iterations suffice. The cost in time on the device is the price of any such step: \benchTenSecRefMin{} to \benchTenSecRefMax{} s per hologram on the benchmark's geometries at 10 outer iterations against \benchTenSecMsevenMin{} to \benchTenSecMsevenMax{} s for phase-fixed WGS, and \benchSecGeoOneBOne{} s on the square lattice at the 100-iteration row against \benchSecGeoOneMseven{} s for phase-fixed WGS; on the acceptance task \secBXMedian{} s against \secMXSevenHundred{} s. It belongs to a hologram that is computed offline and displayed for the length of an experiment, not to one that is recomputed in the loop. Where the device is the benchmark's, with quantization, a fill factor and flicker, the reading at equal time is a gain on the square lattice, the ring and the random set and a tie elsewhere; a user whose geometry is the quasi-periodic set or a large set should expect no gain at equal time. The safeguard did not act on the lattice task and is a cheap guard for geometries where the last iterate loses light. Neither choice of loss needs recommending: the paper measures no advantage for the efficiency-level loss.

The refinement does not replace camera feedback: the paper measures hologram quality in a model, and a flat hologram on the device model is a better start for feedback, not a substitute for it. What would change these conclusions is a measured device model: with a different kernel, a different quantizer or a measured flicker the factors would change, and the sensitivity of the refinement to a mismatched kernel was not tested. The recommendation to keep the weight update off holds on the acceptance task and on the full device and is mixed on the benchmark lattices, where on the ring the update gave the lower median spread (Table~\ref{tab:ablb1geo}), so the option is worth trying on that kind of geometry.
